\chapter{Support Vector Machines and Supervised Synthesis}
\label{ch:svm}

\chapterguide%
  {Linear classifiers from \chapref{ch:classification}, model-capacity ideas from \chapref{ch:training}, and dot products/norms from \chapref{ch:linear-algebra}.}
  {explain margins, kernels, scaling, and SVM tuning; compare supervised model families.}

A support vector machine (SVM) chooses a separating boundary with a wide \emph{margin}. The \term{support vectors}, on or inside that margin, determine the fitted boundary.

\section{The widest street}

\begin{mlfigure}
\begin{tikzpicture}[scale=0.85, font=\small]
% class + (filled)
\foreach \p in {(1,1),(2,1.2),(1,2.6),(0.8,1.8),(2.4,0.8)}
  \fill \p circle (2.4pt);
% class - (open)
\foreach \p in {(4,4),(5,3.5),(3.6,4.9),(4.6,4.7),(5.2,2.9)}
  \draw[thick, fill=gray!30] \p circle (2.4pt);
% support vectors
\fill (1.5,2.5) circle (2.4pt);
\draw[thick, fill=gray!30] (2.5,3.5) circle (2.4pt);
\draw[thick, fill=gray!30] (4,2) circle (2.4pt);
\draw (1.5,2.5) circle (5.5pt);
\draw (2.5,3.5) circle (5.5pt);
\draw (4,2) circle (5.5pt);
% lines
\draw[thick] (0,5) -- (5,0)
  node[pos=0.06, right=2pt] {$\bm{w}^{\mathsf T}\bm{x}+b=0$};
\draw[dashed] (0,4) -- (4,0);
\draw[dashed] (0,6) -- (6,0);
% margin arrow
\draw[{Stealth[length=2mm]}-{Stealth[length=2mm]}, thick]
  (3.1,0.9) -- (4.1,1.9) node[midway, below right=-1pt] {margin};
\node[align=left] at (5.35,5.3) {circled points =\\support vectors};
\end{tikzpicture}
\end{mlfigure}

\begin{intuition}
A hard-margin SVM builds the widest street between two classes. Points touching its curbs determine the street; distant points do not.
\end{intuition}

\section{The hard-margin objective}

For $y_i\in\{-1,+1\}$, require correctly signed linear scores of magnitude at least one:
\[
  \min_{\bm{w},b}\ \frac{1}{2}\|\bm{w}\|_2^2
  \quad\text{subject to}\quad
  y_i(\bm{w}^{\mathsf T}\bm{x}_i+b)\geq 1 .
\]
The constraint requires the correct sign and score magnitude at least~1. The margin width is $2/\|\bm w\|_2$, so minimizing the norm maximizes that width.

\section{The soft margin: allowing exceptions}

Real classes overlap, so the SVM used in practice permits violations and charges for them:
\[
  \min_{\bm{w},b}\
  \frac{1}{2}\|\bm{w}\|_2^2
  +C\sum_{i=1}^{n}
  \max\!\left(0,\,1-y_i(\bm{w}^{\mathsf T}\bm{x}_i+b)\right).
\]
\term{Hinge loss} is zero beyond the margin and linear within it or on the wrong side. $C$ balances margin width against violations:

\begin{itemize}
  \item \term{Large $C$}: violations cost more; weaker regularization can overfit.
  \item \term{Small $C$}: more violations are tolerated; regularization is stronger.
\end{itemize}

Both SVMs and regularized logistic regression combine loss with an $L_2$ penalty. Hinge loss becomes zero beyond the margin; cross-entropy continues rewarding confidence.

\section{Curved boundaries: mapping up}

Some datasets cannot be separated by any straight line --- but they can be after the features are transformed.

\begin{workedexample}
Let $x\in\{-2,2\}$ be class A and $x\in\{-1,0,1\}$ class B. No single threshold separates them. Map to $\phi(x)=(x,x^2)$:
\begin{mltable}
\begin{tabular}{@{}lccccc@{}}
\toprule
\rowcolor{MLPaleBlue!92}
$x$ & $-2$ & $-1$ & $0$ & $1$ & $2$ \\
\midrule
$x^2$ & $4$ & $1$ & $0$ & $1$ & $4$ \\
class & A & B & B & B & A \\
\bottomrule
\end{tabular}
\end{mltable}
Class A has height 4 and class B at most 1. The line $x^2=2.5$ separates them, corresponding to two thresholds $|x|=\sqrt{2.5}$ in the original space.
\end{workedexample}

\section{The kernel trick}

Rich feature maps can be expensive or infinite-dimensional. SVM equations use them only through inner products; a \term{kernel function} computes these directly:
\[
  k(\bm{x},\bm{x}')=\phi(\bm{x})^{\mathsf T}\phi(\bm{x}'),
\]
so the mapped features never need to be constructed at all --- this shortcut is the \term{kernel trick}. Standard choices:

\begin{mltable}
\begin{tabular}{@{}lll@{}}
\toprule
\rowcolor{MLPaleBlue!92}
Kernel & Formula & Boundary it produces \\
\midrule
Linear & $\bm{x}^{\mathsf T}\bm{x}'$ & Straight; fast on huge, sparse data \\
Polynomial & ${(\gamma\,\bm{x}^{\mathsf T}\bm{x}'+r)}^d$ & Curved, degree-$d$ interactions \\
RBF (Gaussian) & $\exp\!\left(-\gamma\|\bm{x}-\bm{x}'\|_2^2\right)$ & Smooth, local, very flexible \\
\bottomrule
\end{tabular}
\end{mltable}

The \term{RBF kernel} measures similarity that decays with distance. Small $\gamma$ gives broad influence and smooth boundaries; large $\gamma$ gives local, flexible boundaries. Tune $C$ and $\gamma$ together on logarithmic grids.

\begin{intuition}
RBF support vectors act as landmarks with distance-decaying influence. The kernel evaluates their similarity without explicitly constructing the feature map.
\end{intuition}

% === VISUAL ADDITION: svm-linear-vs-rbf ===
\subsection{Linear versus RBF boundaries, drawn}

A feature map can make a nonlinear separation linear in the transformed space.

\begin{mlfigure}
\begin{tikzpicture}[font=\scriptsize]
\begin{scope}
  \draw[draw=MLMuted!70,fill=white] (0,0) rectangle (4.7,3.8);
  \node[font=\small\bfseries\color{MLNavy}] at (2.35,4.2) {Linear kernel};
  \foreach \p in {(2.0,1.7),(2.4,1.9),(2.1,2.3),(2.7,2.2),(2.55,1.55)}
    \fill[MLBlue] \p circle (2.6pt);
  \foreach \p in {(0.65,0.65),(1.0,3.05),(3.75,0.7),(4.0,3.0),(0.55,2.0),(4.15,1.85)}
    \draw[MLOrange,very thick] \p circle (2.9pt);
  \draw[MLNavy,very thick] (0.35,2.9) -- (4.35,0.85);
  \node[align=center,text=MLRed] at (2.35,-0.48)
    {one straight line cannot\\put the center in its own class};
\end{scope}

\begin{scope}[xshift=6.0cm]
  \draw[draw=MLMuted!70,fill=white] (0,0) rectangle (4.7,3.8);
  \node[font=\small\bfseries\color{MLNavy}] at (2.35,4.2) {RBF kernel};
  \foreach \p in {(2.0,1.7),(2.4,1.9),(2.1,2.3),(2.7,2.2),(2.55,1.55)}
    \fill[MLBlue] \p circle (2.6pt);
  \foreach \p in {(0.65,0.65),(1.0,3.05),(3.75,0.7),(4.0,3.0),(0.55,2.0),(4.15,1.85)}
    \draw[MLOrange,very thick] \p circle (2.9pt);
  \draw[MLTeal,very thick] (2.35,1.95) ellipse (1.15 and 0.92);
  \node[align=center,text=MLGreen] at (2.35,-0.48)
    {a smooth curved boundary\\can isolate the local region};
\end{scope}

\node[font=\scriptsize\color{MLBlue}] at (4.9,4.18) {$\bullet$ class A};
\node[font=\scriptsize\color{MLOrange}] at (6.65,4.18) {$\circ$ class B};
\end{tikzpicture}
\end{mlfigure}

\begin{keyidea}
A kernel changes the representation during fitting; it does not bend an already-trained line afterward.
\end{keyidea}
% === END VISUAL ADDITION ===

\section{What $C$ and $\gamma$ actually do}
\label{sec:svm-controls}

\begin{mlfigure}
\begin{tikzpicture}[scale=0.82, every node/.style={font=\scriptsize}]
\begin{scope}
\draw[draw=MLMuted,fill=white] (0,0) rectangle (3.5,2.8);
\draw[MLNavy,very thick] (0.3,0.2) -- (3.2,2.6);
\draw[MLNavy,dashed] (0.05,0.75) -- (2.7,2.95);
\draw[MLNavy,dashed] (0.75,-0.15) -- (3.5,2.15);
\fill[MLBlue!12] (0.05,0.75) -- (2.7,2.95) -- (3.5,2.15) -- (0.75,-0.15) -- cycle;
\draw[MLNavy,very thick] (0.3,0.2) -- (3.2,2.6);
\foreach \p in {(0.5,1.9),(1.0,2.4),(0.35,1.35),(1.5,2.65)} \fill[MLBlue] \p circle (2.2pt);
\foreach \p in {(2.4,0.5),(2.9,1.0),(3.2,0.4),(1.9,0.25)} \fill[MLOrange] \p circle (2.2pt);
\fill[MLOrange] (1.65,1.6) circle (2.2pt);
\node[font=\small\bfseries\color{MLNavy}] at (1.75,3.15) {small $C$: wide street};
\node[color=MLMuted,align=center] at (1.75,-0.55) {tolerates the stray point\\simpler, may underfit};
\end{scope}
\begin{scope}[xshift=4.6cm]
\draw[draw=MLMuted,fill=white] (0,0) rectangle (3.5,2.8);
\fill[MLRed!10] (0.35,0.95) -- (2.15,2.9) -- (2.6,2.5) -- (0.8,0.55) -- cycle;
\draw[MLRed,very thick,rounded corners=2pt] (0.2,0.5) -- (1.2,1.3) -- (1.45,1.85) -- (2.0,2.2) -- (2.6,2.8);
\foreach \p in {(0.5,1.9),(1.0,2.4),(0.35,1.35),(1.5,2.65)} \fill[MLBlue] \p circle (2.2pt);
\foreach \p in {(2.4,0.5),(2.9,1.0),(3.2,0.4),(1.9,0.25)} \fill[MLOrange] \p circle (2.2pt);
\fill[MLOrange] (1.65,1.6) circle (2.2pt);
\draw[MLRed,very thick] (1.4,1.45) circle (0.32);
\node[font=\small\bfseries\color{MLNavy}] at (1.75,3.15) {large $C$: narrow street};
\node[color=MLMuted,align=center] at (1.75,-0.55) {bends to capture it\\complex, may overfit};
\end{scope}
\end{tikzpicture}
\end{mlfigure}

\begin{mltable}
\begin{tabularx}{\linewidth}{@{}>{\bfseries\leavevmode\color{MLNavy}}p{0.10\linewidth}p{0.27\linewidth}X@{}}
\toprule
\rowcolor{MLPaleBlue!92}
\normalfont\bfseries\leavevmode\color{MLNavy}Dial & \normalfont\bfseries\leavevmode\color{MLNavy}Plain meaning & \normalfont\bfseries\leavevmode\color{MLNavy}Turn it up when \\
\midrule
$C$ & How badly you want to avoid mistakes on the training data & The model underfits; the boundary is too simple \\
$\gamma$ & How far a single training point's influence reaches & The boundary needs to be more local and wiggly \\
\bottomrule
\end{tabularx}
\end{mltable}

\begin{analogy}
$C$ is a strictness setting for a teacher marking homework. A low $C$ teacher accepts that a few students will get things wrong and teaches a general rule. A high $C$ teacher insists that every single student's answer be accommodated, and ends up teaching a rule so contorted that nobody can apply it to a new question.

$\gamma$ is the radius of each training point's voice. Low $\gamma$ means every point is heard across the whole room, producing a smooth global boundary. High $\gamma$ means each point is only heard by its immediate neighbors, producing islands around individual examples.
\end{analogy}

\section{Multiclass SVMs and support vector regression}

Multiclass reductions include \term{one-vs-rest} ($K$ classifiers) and \term{one-vs-one} (one per class pair, combined by voting).

\term{Support vector regression (SVR)} recycles the margin idea for numeric labels: build a tube of half-width $\varepsilon$ around the prediction and charge nothing for points inside it, using the $\varepsilon$-insensitive loss $\max(0,\,|y-f(\bm{x})|-\varepsilon)$. Only points on or outside the tube become support vectors, giving a sparse, outlier-tolerant regressor with the full kernel toolbox available.

\section{Practical notes}

\begin{itemize}
  \item \term{Scale the features.} Apply the distance-geometry rule from \secref{sec:scaling}.
  \item \term{Tune $C$ (and $\gamma$)} by cross-validation over logarithmic grids such as $10^{-2},\ldots,10^{3}$.
  \item \term{Scale limits:} kernel SVMs suit small-to-medium datasets; linear SVMs suit high-dimensional text. Large datasets may favor linear models or ensembles.
  \item \term{Probabilities are not native.} SVMs output scores and margins; if calibrated probabilities are needed, apply and validate a calibration step (\chapref{ch:evaluation}).
\end{itemize}

\section{Choosing a supervised model}
\label{sec:supervised-model-choice}

\begin{mltable}
\begin{tabular}{@{}p{0.19\linewidth}p{0.36\linewidth}p{0.36\linewidth}@{}}
\toprule
\rowcolor{MLPaleBlue!92}
Model & Useful strengths & Important limitations \\
\midrule
Linear regression & Interpretable, fast, the standard first baseline & Linear relationships only; sensitive to outliers \\
Ridge / lasso & Controls overfitting; lasso selects features & Still linear; needs feature scaling \\
Logistic regression & Strong probability baseline, convex training, interpretable odds & Linear boundary; calibration still needs checking \\
LDA / QDA & Fast, stable on small data & Gaussian class assumptions \\
$k$-NN & No training, flexible boundaries & Slow at prediction; scale- and dimension-sensitive \\
Naive Bayes & Very fast, strong text baseline & Independence and calibration caveats (\secref{sec:naive-bayes}) \\
Decision tree & Readable rules, no scaling, mixed data & Unstable; overfits without pruning \\
Random forest & Robust, accurate, minimal tuning & Larger, slower, less transparent \\
Gradient boosting & Often a very strong tabular benchmark & Needs careful tuning and early stopping \\
SVM & Margins and kernels; strong in high dimensions & Scaling required; slow on very large data \\
\bottomrule
\end{tabular}
\end{mltable}

Start with regularized linear models; add complexity only for validated gains. In Lab~3, a tuned SVM reaches test AUC 0.997 versus 0.986--0.992 for simpler KNN/naive Bayes fits.

\begin{keyidea}
SVMs illustrate two reusable ideas: wide margins control complexity, and kernels allow linear methods to express nonlinear boundaries.
\end{keyidea}

\begin{modelcard}{Support Vector Machines - Quick Guide}
\begin{tabularx}{\linewidth}{@{}>{\bfseries\leavevmode\color{MLNavy}}p{0.24\linewidth}X@{}}
Best when & The dataset is small to medium sized, features are informative, and a strong margin-based boundary is useful. \cr
Preprocessing & Standardize numerical features; encode categories; remove clearly irrelevant features when justified by the validated workflow in \secref{sec:feature-selection}. \cr
Main controls & $C$ for margin violations; kernel choice; $\gamma$ for RBF-kernel reach; $\varepsilon$ for SVR. \cr
Strengths & Effective in high-dimensional spaces and flexible through kernels. \cr
Watch for & Slow training on very large datasets, sensitive scale, difficult probability interpretation, and expensive kernel tuning. \cr
\end{tabularx}
\end{modelcard}

\labsection{3}{Supervised learning II: model families}{sec:lab03-supervised-models}

\begin{labproject}{Neighbours, naive Bayes, trees, ensembles, and support vector machines on one workflow}
\textbf{Goal.} Compare model families under Lab~2's fixed validation workflow; inspect boundaries, tree depth, ensemble gains, and SVM tuning.

\medskip\noindent\textbf{Setup.}\par
\begin{lstlisting}[style=mlpython]
from pathlib import Path

import joblib
import matplotlib.pyplot as plt
import numpy as np
import pandas as pd
from sklearn.compose import ColumnTransformer
from sklearn.ensemble import GradientBoostingRegressor, RandomForestRegressor
from sklearn.impute import SimpleImputer
from sklearn.inspection import DecisionBoundaryDisplay, permutation_importance
from sklearn.linear_model import LogisticRegression
from sklearn.metrics import accuracy_score, mean_absolute_error, mean_squared_error, roc_auc_score
from sklearn.model_selection import (
    GridSearchCV, KFold, StratifiedKFold, cross_val_score, train_test_split,
)
from sklearn.naive_bayes import GaussianNB
from sklearn.neighbors import KNeighborsClassifier
from sklearn.pipeline import Pipeline
from sklearn.preprocessing import OneHotEncoder, StandardScaler
from sklearn.svm import SVC
from sklearn.tree import DecisionTreeClassifier, export_text, plot_tree

PROJECT_ROOT = Path.cwd().parent
DATA = PROJECT_ROOT / "all_datasets"
MODELS = PROJECT_ROOT / "models"
MODELS.mkdir(exist_ok=True)
\end{lstlisting}

\medskip\noindent\textbf{Part A --- KNN versus Gaussian naive Bayes.}\par

KNN votes among standardized neighbors; Gaussian naive Bayes combines per-feature Gaussian likelihoods and class priors.

\begin{lstlisting}[style=mlpython]
cancer = pd.read_csv(DATA / "breast_cancer_dataset.csv")

X_c = cancer.drop(columns=["id", "diagnosis", "Unnamed: 32"], errors="ignore")
y_c = (cancer["diagnosis"] == "M").astype(int)

Xc_train, Xc_valid, yc_train, yc_valid = train_test_split(
    X_c, y_c, test_size=0.25, random_state=42, stratify=y_c
)
\end{lstlisting}

$k$ is a hyperparameter, so it is chosen by cross-validation inside the training partition. Small $k$ follows every local quirk (low bias, high variance); large $k$ smooths everything (high bias, low variance).

\begin{lstlisting}[style=mlpython]
cv = StratifiedKFold(n_splits=5, shuffle=True, random_state=42)
k_values = [1, 3, 5, 7, 11, 21]
cv_scores = {}

for k in k_values:
    candidate = Pipeline([
        ("scale", StandardScaler()),
        ("model", KNeighborsClassifier(n_neighbors=k, weights="distance")),
    ])
    scores = cross_val_score(candidate, Xc_train, yc_train, cv=cv, scoring="roc_auc")
    cv_scores[k] = scores.mean()
    print(f"k={k:2d}  cross-validated ROC AUC = {scores.mean():.3f}")

best_k = max(cv_scores, key=cv_scores.get)
print("selected k:", best_k)

plt.figure(figsize=(5, 3.2))
plt.plot(k_values, [cv_scores[k] for k in k_values], marker="o")
plt.axvline(best_k, linestyle="--", color="grey")
plt.xlabel("k (number of neighbours)")
plt.ylabel("Cross-validated ROC AUC")
plt.title("Choosing k on training data only")
plt.tight_layout()
plt.show()
\end{lstlisting}

\textbf{Inspect.} Validation improves beyond $k=1$ and is similar for $k=7$--21.

With $k$ locked, both models are compared on the validation split.

\begin{lstlisting}[style=mlpython]
models = {
    "KNN": Pipeline([
        ("scale", StandardScaler()),
        ("model", KNeighborsClassifier(n_neighbors=best_k, weights="distance")),
    ]),
    "GaussianNB": Pipeline([
        ("scale", StandardScaler()),
        ("model", GaussianNB()),
    ]),
}

for name, model in models.items():
    model.fit(Xc_train, yc_train)
    pred = model.predict(Xc_valid)
    prob_c = model.predict_proba(Xc_valid)[:, 1]
    print(f"{name:10s} accuracy={accuracy_score(yc_valid, pred):.3f} "
          f"roc_auc={roc_auc_score(yc_valid, prob_c):.3f}")
\end{lstlisting}

Compare two-feature boundaries on tumor radius and texture, adding logistic regression as a reference.

\begin{lstlisting}[style=mlpython]
two = ["radius_mean", "texture_mean"]
fig, axes = plt.subplots(1, 3, figsize=(12, 3.8))
for ax, (name, estimator) in zip(axes, [
    ("Logistic regression", LogisticRegression(max_iter=2000)),
    (f"KNN (k={best_k})", KNeighborsClassifier(n_neighbors=best_k, weights="distance")),
    ("Gaussian naive Bayes", GaussianNB()),
]):
    small = Pipeline([("scale", StandardScaler()), ("model", estimator)])
    small.fit(Xc_train[two], yc_train)
    DecisionBoundaryDisplay.from_estimator(small, Xc_valid[two], ax=ax, alpha=0.3,
                                           response_method="predict", cmap="coolwarm")
    ax.scatter(Xc_valid[two[0]], Xc_valid[two[1]], c=yc_valid, cmap="coolwarm",
               s=12, edgecolor="k", linewidth=0.3)
    ax.set_title(f"{name}: {accuracy_score(yc_valid, small.predict(Xc_valid[two])):.2f} accuracy")
    ax.set_xlabel("radius_mean")
    ax.set_ylabel("texture_mean")
plt.tight_layout()
plt.show()
\end{lstlisting}

\textbf{Inspect.} Logistic regression is linear; Gaussian naive Bayes curves; KNN follows local votes.

\medskip\noindent\textbf{Part B --- One readable decision tree, and what happens as it grows.}\par

Trees split on value order, so numerical preprocessing imputes missingness without scaling.

\begin{lstlisting}[style=mlpython]
titanic = pd.read_csv(DATA / "titanic_dataset.csv")

features = ["Pclass", "Sex", "Age", "SibSp", "Parch", "Fare", "Embarked"]
t_num = ["Age", "SibSp", "Parch", "Fare"]
t_cat = ["Pclass", "Sex", "Embarked"]
X_t = titanic[features]
y_t = titanic["Survived"]

tree_preprocess = ColumnTransformer([
    ("num", SimpleImputer(strategy="median"), t_num),
    ("cat", Pipeline([
        ("impute", SimpleImputer(strategy="most_frequent")),
        ("onehot", OneHotEncoder(handle_unknown="ignore")),
    ]), t_cat),
])

Xt_train, Xt_valid, yt_train, yt_valid = train_test_split(
    X_t, y_t, test_size=0.25, random_state=42, stratify=y_t
)

tree = Pipeline([
    ("preprocess", tree_preprocess),
    ("tree", DecisionTreeClassifier(max_depth=4, min_samples_leaf=12, random_state=42)),
])
tree.fit(Xt_train, yt_train)
print("training accuracy:", round(accuracy_score(yt_train, tree.predict(Xt_train)), 3))
print("validation accuracy:", round(accuracy_score(yt_valid, tree.predict(Xt_valid)), 3))
\end{lstlisting}

Print and plot the tree; inspect split questions, sample counts, and classes.

\begin{lstlisting}[style=mlpython]
feature_names = list(tree.named_steps["preprocess"].get_feature_names_out())
print(export_text(tree.named_steps["tree"], feature_names=feature_names, max_depth=3))

plt.figure(figsize=(12, 6))
plot_tree(tree.named_steps["tree"], feature_names=feature_names,
          class_names=["did not survive", "survived"], max_depth=2, filled=True, fontsize=8)
plt.title("Top of the fitted Titanic decision tree")
plt.tight_layout()
plt.show()
\end{lstlisting}

\textbf{Inspect.} The tree first splits on sex, then age for men and class for women.

Fit each depth and compare training with validation accuracy.

\begin{lstlisting}[style=mlpython]
depths = list(range(1, 13))
train_acc, valid_acc = [], []
for depth in depths:
    t = Pipeline([
        ("preprocess", tree_preprocess),
        ("tree", DecisionTreeClassifier(max_depth=depth, random_state=42)),
    ]).fit(Xt_train, yt_train)
    train_acc.append(accuracy_score(yt_train, t.predict(Xt_train)))
    valid_acc.append(accuracy_score(yt_valid, t.predict(Xt_valid)))

plt.figure(figsize=(6, 3.6))
plt.plot(depths, train_acc, marker="o", label="training")
plt.plot(depths, valid_acc, marker="o", label="validation")
plt.xlabel("max_depth")
plt.ylabel("Accuracy")
plt.title("Deeper trees memorise; validation accuracy stops improving")
plt.legend()
plt.tight_layout()
plt.show()
\end{lstlisting}

\textbf{Inspect.} Training accuracy rises with depth while validation peaks earlier, exposing overfitting.

\medskip\noindent\textbf{Part C --- Random forest versus gradient boosting.}\par

Model $\log(1+\text{price})$ from Ames Housing's mixed-type features. Compare random forest and gradient boosting on identical three-fold splits.

\begin{lstlisting}[style=mlpython]
ames = pd.read_csv(DATA / "ames_housing_dataset.csv")
X_a = ames.drop(columns=["SalePrice", "Order", "PID"], errors="ignore")   # identifiers out
y_a = np.log1p(ames["SalePrice"])

ames_preprocess = ColumnTransformer([
    ("num", SimpleImputer(strategy="median"), X_a.select_dtypes(include="number").columns),
    ("cat", Pipeline([
        ("impute", SimpleImputer(strategy="most_frequent")),
        ("onehot", OneHotEncoder(handle_unknown="ignore", sparse_output=False)),
    ]), X_a.select_dtypes(exclude="number").columns),
])
ensembles = {
    "random_forest": RandomForestRegressor(
        n_estimators=150, max_features=0.8, min_samples_leaf=2, random_state=42, n_jobs=-1
    ),
    "gradient_boosting": GradientBoostingRegressor(
        n_estimators=150, learning_rate=0.05, max_depth=3, loss="huber", random_state=42
    ),
}

cv3 = KFold(n_splits=3, shuffle=True, random_state=42)
cv_rmse = {}
for name, estimator in ensembles.items():
    pipe = Pipeline([("preprocess", ames_preprocess), ("model", estimator)])
    neg_mse = cross_val_score(pipe, X_a, y_a, cv=cv3, scoring="neg_mean_squared_error")
    cv_rmse[name] = np.sqrt(-neg_mse)
    print(f"{name:18s} CV log-RMSE = {cv_rmse[name].mean():.4f} +/- {cv_rmse[name].std():.4f}")

plt.figure(figsize=(5, 3.4))
for i, (name, values) in enumerate(cv_rmse.items()):
    plt.scatter([i] * len(values), values, s=60, label="fold scores" if i == 0 else None)
    plt.hlines(values.mean(), i - 0.2, i + 0.2, color="black", label="mean" if i == 0 else None)
plt.xticks(range(len(cv_rmse)), list(cv_rmse))
plt.ylabel("log-RMSE per fold (lower is better)")
plt.title("Fold-by-fold: is the gap bigger than the spread?")
plt.legend()
plt.tight_layout()
plt.show()
\end{lstlisting}

\textbf{Inspect.} Boosting's observed gain is modest relative to fold variation. Weigh it against tuning and runtime.

Repeat on eleven wine measurements. Permute validation columns to measure the fitted forest’s reliance on each.

\begin{lstlisting}[style=mlpython]
wine = pd.read_csv(DATA / "red_wine_quality_dataset.csv")
X_w = wine.drop(columns="quality")
y_w = wine["quality"]
Xw_train, Xw_valid, yw_train, yw_valid = train_test_split(X_w, y_w, test_size=0.25, random_state=42)

wine_models = {
    "random_forest": RandomForestRegressor(n_estimators=400, min_samples_leaf=2,
                                           random_state=42, n_jobs=-1),
    "gradient_boosting": GradientBoostingRegressor(random_state=42),
}
wine_pred = {}
for name, model in wine_models.items():
    model.fit(Xw_train, yw_train)
    wine_pred[name] = model.predict(Xw_valid)
    print(f"{name:18s} MAE = {mean_absolute_error(yw_valid, wine_pred[name]):.3f}  "
          f"RMSE = {np.sqrt(mean_squared_error(yw_valid, wine_pred[name])):.3f}")

importance = permutation_importance(
    wine_models["random_forest"], Xw_valid, yw_valid, n_repeats=10, random_state=42, n_jobs=-1
)
ranked = pd.Series(importance.importances_mean, index=X_w.columns).sort_values()
spread = importance.importances_std[np.argsort(importance.importances_mean)]

fig, axes = plt.subplots(1, 2, figsize=(11, 3.8))
axes[0].barh(ranked.index, ranked.values, xerr=spread)
axes[0].set_xlabel("Increase in error when the column is shuffled")
axes[0].set_title("Permutation importance (random forest)")
axes[1].scatter(yw_valid, wine_pred["gradient_boosting"], alpha=0.5)
axes[1].plot([3, 8], [3, 8], linestyle="--", color="grey")
axes[1].set_xlabel("Observed quality")
axes[1].set_ylabel("Predicted quality")
axes[1].set_title("Gradient boosting pulls extremes to the middle")
plt.tight_layout()
plt.show()
\end{lstlisting}

\textbf{Inspect.} Alcohol and sulphates rank highly; predictions compress rare extreme ratings. Forests beat boosting here, unlike Ames.

\medskip\noindent\textbf{Part D --- Support vector machines: two dials and a kernel.}\par

Fit SVM scaling within each training fold. $C$ controls violation penalties; the kernel controls boundary shape; RBF $\gamma$ controls locality.

\begin{lstlisting}[style=mlpython]
X_train, X_test, y_train, y_test = train_test_split(
    X_c, y_c, test_size=0.25, random_state=42, stratify=y_c
)

svm_pipe = Pipeline([("scale", StandardScaler()), ("model", SVC())])
search = GridSearchCV(
    svm_pipe,
    param_grid={
        "model__kernel": ["linear", "rbf"],
        "model__C": [0.1, 1, 10, 100],
        "model__gamma": ["scale", 0.01, 0.1],
    },
    cv=StratifiedKFold(5, shuffle=True, random_state=42),
    scoring="roc_auc",
    n_jobs=-1,
)
search.fit(X_train, y_train)

score = search.decision_function(X_test)
pred = search.predict(X_test)
print("best parameters:", search.best_params_)
print("cross-validated ROC AUC:", round(search.best_score_, 3))
print("test accuracy:", round(accuracy_score(y_test, pred), 3))
print("test ROC AUC:", round(roc_auc_score(y_test, score), 3))
\end{lstlisting}

Plot cross-validated AUC for each $(C,\gamma)$ pair and compare two-feature kernel boundaries.

\begin{lstlisting}[style=mlpython]
results = pd.DataFrame(search.cv_results_)
rbf = results[results["param_model__kernel"] == "rbf"].copy()
rbf["gamma"] = rbf["param_model__gamma"].astype(str)
grid = rbf.pivot(index="param_model__C", columns="gamma", values="mean_test_score")

fig, axes = plt.subplots(1, 3, figsize=(13, 3.8))
im = axes[0].imshow(grid.values, cmap="viridis", aspect="auto")
axes[0].set_xticks(range(len(grid.columns)), grid.columns)
axes[0].set_yticks(range(len(grid.index)), grid.index)
axes[0].set_xlabel("gamma")
axes[0].set_ylabel("C")
for i in range(grid.shape[0]):
    for j in range(grid.shape[1]):
        axes[0].text(j, i, f"{grid.values[i, j]:.3f}", ha="center", va="center",
                     color="white", fontsize=8)
axes[0].set_title("RBF kernel: CV ROC AUC by C and gamma")
for ax, (title, estimator) in zip(axes[1:], [
    ("Linear kernel", SVC(kernel="linear", C=1)),
    ("RBF kernel (gamma = 1)", SVC(kernel="rbf", C=1, gamma=1.0)),
]):
    small = Pipeline([("scale", StandardScaler()), ("model", estimator)]).fit(X_train[two], y_train)
    DecisionBoundaryDisplay.from_estimator(small, X_test[two], ax=ax, alpha=0.3,
                                           response_method="predict", cmap="coolwarm")
    ax.scatter(X_test[two[0]], X_test[two[1]], c=y_test, cmap="coolwarm",
               s=12, edgecolor="k", linewidth=0.3)
    ax.set_title(f"{title}: {accuracy_score(y_test, small.predict(X_test[two])):.2f} accuracy")
    ax.set_xlabel("radius_mean")
    ax.set_ylabel("texture_mean")
plt.tight_layout()
plt.show()
\end{lstlisting}

\textbf{Inspect.} Most scores are similar. RBF allows curved boundaries, but these scaled classes gain little over a linear model.

Save, reload, and compare predictions with the in-memory model before sharing the artifact.

\begin{lstlisting}[style=mlpython]
model_path = MODELS / "breast_cancer_svm.joblib"
joblib.dump(search.best_estimator_, model_path)
loaded = joblib.load(model_path)
same = np.array_equal(loaded.predict(X_test), pred)
print("file size (KB):", round(model_path.stat().st_size / 1024, 1))
print("reloaded model reproduces the test predictions:", same)
\end{lstlisting}

\begin{tryit}
Grow the Titanic tree with \texttt{max\_depth=None} and \texttt{min\_samples\_leaf=1}, and compare training and validation accuracy with the depth-4 tree. Then refit the SVM search with only the linear kernel and compare its test ROC AUC with the RBF result.
\end{tryit}
\end{labproject}


\begin{quickcheck}
\begin{enumerate}
  \item What is the geometric meaning of a larger separating margin in an SVM?
  \item How do \(C\) and, for an RBF kernel, \(\gamma\) change the flexibility of the fitted decision boundary?
  \item Why should scaling and SVM hyperparameter selection both happen inside the validation procedure?
\end{enumerate}
\end{quickcheck}